% Werte und Aufgabenstellung: aufgaben/Blatt07.pdf, S. 1.

\subsection*{Gegeben und Modell}
Der Draht hat die ungedehnte Länge $l$, den überall gleichen Querschnitt $A$,
die homogene Dichte $\rho$ im ungedehnten Zustand und den Elastizitätsmodul
$E$. Gesucht sind die Verlängerung durch das Eigengewicht und die kritische
Länge für Stahl. Für die Zahlenrechnung verwenden wir
$g=\qty{9.81}{\metre\per\second\squared}$; dieser Näherungswert steht nicht
auf dem Aufgabenblatt.

Wir betrachten den Draht im statischen Gleichgewicht in einem konstanten
Schwerefeld, ohne zusätzliche Last. Material und Querschnitt sind homogen;
die Querkontraktion wird wie vorgegeben vernachlässigt. Die lokale Dehnung
folgt dem linearen Hooke'schen Gesetz. Für die Verlängerung am Bruch nehmen
wir zusätzlich an, dass dieses Gesetz bis zur angegebenen Grenzspannung gilt.
Bei realem Stahl kann vorher plastische Verformung auftreten.

\subsection*{a) Verlängerung durch das Eigengewicht}
Wir messen die \emph{ungedehnte Materialkoordinate} $x$ vom oberen Ende nach
unten. Ein Schnitt bei $x$ trägt den darunter befindlichen Draht der
ungedehnten Länge $l-x$. Weil dessen Masse $\rho A(l-x)$ beträgt, ist die
Zugkraft im statischen Gleichgewicht
\begin{equation}
  F(x)=\rho A g(l-x), \qquad
  \sigma(x)=\frac{F(x)}{A}=\rho g(l-x).
\end{equation}
Die Kraft ist also entlang des Drahts nicht konstant. Die auf einem
ursprünglich $\mathrm dx$ langen Stück lokal entstehende Längenänderung folgt
aus $\varepsilon(x)=\sigma(x)/E$:
\begin{equation}
  \mathrm d(\Delta l)=\varepsilon(x)\,\mathrm dx
  =\frac{\rho g}{E}(l-x)\,\mathrm dx.
\end{equation}
Aufsummieren über alle \emph{ursprünglichen} Stücke liefert
\begin{align}
  \Delta l
    &=\frac{\rho g}{E}\int_0^l(l-x)\,\mathrm dx
      =\frac{\rho g}{E}\left[lx-\frac{x^2}{2}\right]_0^l
      =\frac{\rho g l^2}{2E}.
\end{align}
\begin{ergebnis}[a) Verlängerung]
\begin{equation*}
  \boxed{\Delta l=\frac{\rho g l^2}{2E}}, \qquad
  L_{\mathrm{gedehnt}}=l+\Delta l.
\end{equation*}
\end{ergebnis}
Der Querschnitt kürzt sich heraus: Ein grösserer Querschnitt trägt zwar mehr
Gewicht, verteilt die Kraft aber auf eine entsprechend grössere Fläche.

\subsection*{b) Grenzlänge des Stahldrahts}
Mit zunehmendem Abstand vom unteren Ende nimmt das getragene Gewicht zu.
Daher ist die Spannung am oberen Ende maximal; dort wird die Grenzspannung
zuerst erreicht. Bei $x=0$ gilt
\begin{equation}
  \sigma_{\mathrm{oben}}=\rho g l.
\end{equation}
An der idealisierten Bruchgrenze setzen wir $\sigma_{\mathrm{oben}}=\sigma_{\max}$
und erhalten für die \emph{ungedehnte} Grenzlänge
\begin{equation}
  l_{\mathrm{krit}}=\frac{\sigma_{\max}}{\rho g}.
\end{equation}
Für Stahl sind gegeben:
\begin{align*}
  \rho&=\qty{7.85}{\gram\per\centi\metre\cubed}
       =\qty{7850}{\kilogram\per\metre\cubed},\\
  E&=\qty{210}{\kilo\newton\per\milli\metre\squared}
    =\qty{210e9}{\pascal},\\
  \sigma_{\max}&=\qty{320}{\newton\per\milli\metre\squared}
    =\qty{320e6}{\pascal}.
\end{align*}
Einsetzen ergibt
\begin{ergebnis}[b) Ungedehnte Grenzlänge]
\begin{equation*}
  l_{\mathrm{krit}}
  =\frac{\qty{320e6}{\pascal}}
         {\qty{7850}{\kilogram\per\metre\cubed}\,
          \qty{9.81}{\metre\per\second\squared}}
  \approx\qty{4.16}{\kilo\metre}.
\end{equation*}
\end{ergebnis}
$E$ geht in diese Grenzlänge nicht ein, wohl aber in die zugehörige
Verlängerung. Unter der zusätzlichen Annahme linearer Elastizität bis zur
Grenze folgt aus Teil a
\begin{equation}
  \Delta l_{\mathrm{krit}}
  =\frac{\sigma_{\max}}{2E}\,l_{\mathrm{krit}}
  \approx\qty{3.17}{\metre}, \qquad
  L_{\mathrm{krit}}=l_{\mathrm{krit}}+\Delta l_{\mathrm{krit}}
  \approx\qty{4158.55}{\metre}.
\end{equation}
Hier bezeichnet $L_{\mathrm{krit}}$ die Länge im gedehnten Zustand unmittelbar
an der idealisierten Grenze. Die letzte Zahl ist nur eine rechnerische
Modellgrösse, keine Aussage über zentimetergenaue reale Bruchlängen.

\begin{kontrolle}[Einheiten und Modellprüfung]
$\rho g l^2/E$ und $\sigma_{\max}/(\rho g)$ haben beide die Einheit Länge.
Für $g\to0$ oder $E\to\infty$ verschwindet die Verlängerung. Die lokale
Dehnung ist unten null und oben maximal; ihr Mittelwert beträgt die Hälfte
des oberen Wertes. An der Grenze ist
$\varepsilon_{\mathrm{oben}}=\sigma_{\max}/E\approx\num{0.00152}$,
die mittlere Dehnung also ungefähr $\num{0.000762}$. Das bestätigt die
kleine Dehnung im Modell. Reale Werte können unter anderem durch plastisches
Verhalten, Querkontraktion und Materialfehler abweichen.
\end{kontrolle}

\begin{merksatz}
Bei veränderlicher Belastung zuerst die lokale Schnittkraft bestimmen. Das
Stoffgesetz gilt lokal; die gesamte Längenänderung entsteht durch Addition
über die Abschnitte in ihrer ursprünglichen Geometrie.
\end{merksatz}
